% =============================================================================
%  02_complexes.tex —— §2 链复形
% =============================================================================
\section{链复形}\label{sec:cx}

\subsection*{拓扑中的链复形}

链复形在代数拓扑里是熟面孔。一个单纯复形 $X$（例如流形的一个三角剖分）给出
两个链复形。取 $C_i$ 为空间中 $i$ 维单形生成的自由阿贝尔群，$C^i$ 为其对偶，
\begin{equation}\label{eq:co}
C_i=\Z\{i\text{-单形}\},\qquad C^i=\Hom(C_i,\Z),
\end{equation}
便得到复形
$$
C\ldot=\ldots\to C_i\xrightarrow{\partial_i}C_{i-1}\to\ldots,
\qquad
C\udot=\ldots\to C^{i-1}\xrightarrow{d^i}C^i\to\ldots.
$$
这里 $\partial$ 是把单形送到其边界的边界算子，$d$ 是它的转置，且
$\partial^2=0=d^2$。于是可以取同调群与上同调群
$$
H_i(X)=\frac{\ker\partial_i}{\operatorname{im}\partial_{i+1}},
\qquad
H^i(X)=\frac{\ker d^{i+1}}{\operatorname{im} d^i}.
$$
单纯映射 $X\to Y$ 诱导相应链复形上的链映射（即与（上）边界算子交换的映射），
从而诱导（上）同调群上的映射。

然而（上）同调 $H_*\ (H^*)$ 所含的信息比复形 $C\ldot$、$C\udot$ 少，
例如 Massey 积就看不见\yizhuRef{4}。存在同调群 $H_*$ 相同、同伦型却不同的空间，
所以 $H_*$ 是个有局限的同伦型不变量。但由 Whitehead 定理，$C\ldot$ 是个非常强的
“不变量”，至少对单连通空间如此（本文只讨论这类空间）。这定理说：单纯复形
$X$、$Y$（都单连通）的底拓扑空间 $|X|$、$|Y|$ 同伦等价，当且仅当存在单纯复形
之间的映射
\begin{equation}\label{eq:wh}
\begin{array}{ccc}
& Z & \\
\swarrow & & \searrow \\
X & & Y
\end{array}
\qquad\Longrightarrow\qquad
\begin{array}{ccc}
& C\ldot^{Z} & \\
\swarrow & & \searrow \\
C\ldot^{X} & & C\ldot^{Y}
\end{array}
\end{equation}
\emph{诱导同调上的同构} $H_*(X)\xleftarrow{\ \sim\ }H_*(Z)\xrightarrow{\ \sim\ }H_*(Y)$
\yizhuRef{5}。（$Z$ 之所以出现，是因为我们可能需要先细分 $X$：取一个 $Z$，使有
单纯逼近 $Z\to X$ 给出同胚 $|Z|\xrightarrow{\sim}|X|$；然后才能对给定的同伦等价
$|X|\to|Y|$ 找到单纯逼近 $Z\to Y$。）

所以 $C\ldot$ 含有的信息已经多到我们不能指望更多，我们愿意把它看成拓扑空间
$|X|$ 的不变量。这件事怎么做稍后再说；眼下可以把它看成三角剖分（单纯复形）的
不变量。

拓扑里还有一个地方，复形优于同调的做法早已是常识，那就是对偶上同调理论的
构造。具体说来，我们\emph{不}把 $H^*$ 定义成 $H_*$ 的对偶，即不取
$$
H^*\ne\Hom(H_*,\Z).
$$
原因在于 $\Hom(\,\cdot\,,\Z)$ 会抹掉挠信息，而且它的平方不是恒等映射
（这个意义上的双重对偶不是原来的 $H_*$）。我们改为在链复形的层面上取
$\Hom(\,\cdot\,,\Z)$，如 \eqref{eq:co} 所示。这样不丢信息，双重对偶就是原来的
复形。

例如把 $\Hom(\,\cdot\,,\Z)$ 作用到复形
$$
\Z\xrightarrow{2}\Z\qquad H_1=0,\ H_0=\Z/2
$$
上，得到对偶复形
$$
\Z\xrightarrow{2}\Z\qquad H^1=\Z/2,\ H^0=0.
$$
可见挠群 $\Z/2$ 在对偶之下并没有丢失（若用 $\Hom(H_*,\Z)$ 就会丢掉），
只是挪到了另一个次数。

至此我们看到了两个例子，它们体现了本次报告的主题：
\begin{center}
\begin{tcolorbox}[colback=black!2!white, colframe=black!60!black, boxrule=0.8pt,
  arc=0pt, width=8.6cm, halign=center, top=2pt, bottom=2pt]
复形好，（上）同调坏\\[2pt]
{\small\itshape Complexes good, (co)homology bad}
\end{tcolorbox}
\end{center}

\begin{yizhu}[Massey 积为什么“看不见”]
同调群只记得“哪条链闭、哪条链是边界”，却不记得链本身。Massey 积是三个上同调类
之间的二次运算，它依赖代表元的选择，并且会因代表元的同调替换而改变；因此它不能
由 $H^*$ 上的运算给出，只能由复形（或它的 $A_\infty$ 结构）提供。这就是作者说
“$H_*$ 所含信息更少”的确切含义。
[来源：Gelfand--Manin 第 IV 章关于 $A_\infty$ 与 Massey 积的讨论；A 级]
\end{yizhu}

\begin{yizhu}[Whitehead 定理里的 $Z$]
底拓扑空间同伦等价，并不保证存在单纯映射 $X\to Y$（同伦等价可以非常“不单纯”）。
标准做法是先用单纯逼近定理：细分 $X$ 得到 $Z$，使 $|Z|\to|X|$ 是同胚，
$Z\to X$ 是单纯逼近；再把给定的同伦等价 $|X|\to|Y|$ 单纯逼近成 $Z\to Y$。
这时两个链映射在 $H_*$ 上都是同构。
[来源：原文此处括注；Whitehead 定理的标准陈述见代数拓扑教材；A 级]
\end{yizhu}

\begin{yizhu}[双重对偶为什么不是恒等]
对有限生成自由 $\Z$-模，$\Hom(\Hom(M,\Z),\Z)\cong M$；但对挠模，例如
$M=\Z/2$，有 $\Hom(\Z/2,\Z)=0$，取两次对偶仍得 $0\ne\Z/2$。所以“
$\Hom(-,\Z)$ 与 $\Hom(H_*,\Z)$ 可交换”这种想法是错的；正确做法正如作者所说，
把对偶放到复形层面：自由复形的双重对偶回到原复形。
[来源：同调代数的标准事实；A 级]
\end{yizhu}

\subsection*{代数几何中的链复形}

在代数几何里考虑链复形，乍看有些奇怪。它出现的主要场合是层的消解
（resolution），下面就来描述。凝聚层常常自然地作为映射的核或余核出现，
也就是作为该映射给出的两项复形的上同调出现。例如设 $D\subset X$ 是光滑代数簇
$X$ 中的除子，它对应一个线丛 $L$ 与一个在 $D$ 上消失的截面 $s\in H^0(L)$
（我们把 $L$ 与其截面层 $L$ 混为一谈）。这就给出标准正合列
\begin{equation}\label{eq:div}
0\to L^{-1}\xrightarrow{\ s\ }\OO_X\to\OO_D\to 0,
\end{equation}
其中 $\OO_X$ 是 $X$ 的结构层（$X$ 上平凡线丛的截面层），$\OO_D$ 是 $D$ 的结构层
推前到 $X$（零延拓；这是集中在 $D$ 上的挠层）。于是 $\OO_D$ 是复形
$\{L^{-1}\xrightarrow{s}\OO_X\}$ 的上同调\yizhuRef{6}。（注意这是层的复形的
上同调，因此仍是一个层；不要与取值在向量空间的层上同调混淆。）

同样，若 $Z\subset X$ 是余维 $r$ 的子簇，且是秩 $r$ 向量丛 $E$ 的正则截面
$s\in H^0(E)$ 的横截零点轨迹，那么正合列（Koszul 复形）
$$
0\to\Lambda^rE^*\to\Lambda^{r-1}E^*\to\ldots\to E^*\to\OO_X,
$$
其中每个箭头是与 $s$ 作内乘\yizhuRef{7}，按定义其余核为 $\OO_Z$。于是 $\OO_Z$
是这个复形的上同调。

这里其实有所得：我们把难缠的挠层 $\OO_D$、$\OO_Z$ 换成了 $X$ 上更驯服的局部自由层
（即向量丛）。一般地可以考察这样的\term{消解}{resolution}：把任意层 $\F$ 换成
在某些意义上“更好”的层组成的复形 $F\udot$，
$$
F\udot\to\F\qquad\text{或}\qquad \F\to F\udot,
$$
然后拿复形 $F\udot$ 代替它那不易驾驭的上同调 $\F$。可以这样想：难看的弯曲的
$\F$ 由较直的 $F$ 拼成，
$$
\F=F^1-F^2+F^3-\ldots,
$$
其中“相减”的含义由消解中各项之间的映射给出，$F^i$ 则是积木：$\F$ 的生成元给出
$F^1$，关系给出 $F^2$，关系之间的关系给出 $F^3$，依此类推\yizhuRef{8}。

当然，消解一般不会自然出现，必须挑一个；这与给拓扑空间挑一个（非典范的）
三角剖分完全类似，我们会在 §\ref{sec:qi} 看到。导出范畴要解决的正是这个问题。
至于层组成的复形 $F\udot$ 在什么意义上“更好”，§\ref{sec:df} 会一般地处理，
下面先看一个简单例子。

\begin{yizhu}[复形的上同调还是层]
初学者容易把“$\OO_D$ 是复形 $\{L^{-1}\to\OO_X\}$ 的上同调”误读成
“$\OO_D$ 是 $D$ 的层上同调”。前者是层的复形逐点取核模余核，结果仍是 $X$ 上的层；
后者 $H^i(D,\OO_D)$ 取值在向量空间，是完全不同的对象。作者在这里专门提醒一句，
是因为后文（§\ref{sec:df}）要用前者的语言重述后者。
[来源：原文此处括注；A 级（原文自述）]
\end{yizhu}

\begin{yizhu}[Koszul 复形的内乘]\label{yizhu:koszul}
设 $E$ 秩为 $r$，$s\in H^0(E)$。内乘 $\iota_s\colon\Lambda^kE^*\to\Lambda^{k-1}E^*$
由 $\iota_s(\xi_1\wedge\ldots\wedge\xi_k)=\sum_j(-1)^{j-1}\xi_j(s)\,
\xi_1\wedge\ldots\widehat{\xi_j}\ldots\wedge\xi_k$ 给出。正则性保证
$\iota_s$ 的余核恰为 $\OO_Z$，且复形在 $Z$ 外正合——这就是“Koszul 复形消解
$\OO_Z$”的内容。
[来源：标准构造，见 Hartshorne《Residues and Duality》第 II 章；A 级]
\end{yizhu}

\begin{yizhu}[$\F=F^1-F^2+F^3-\ldots$ 该怎么读]
这不是等式，而是“把 $\F$ 想象成由 $F^i$ 逐层搭起来”的记法：自由消解中，
$F^1$ 由 $\F$ 的生成元张成，$F^2$ 记录生成元之间的关系，$F^3$ 记录关系之间的
关系。复形的微分把这些积木粘起来；所谓“相减”，就是微分中出现的符号。
[来源：译者按（据原文上下文与自由消解的标准构造）；无外部页码依据]
\end{yizhu}

\subsection*{用层做相交理论}

我们只看最简单的情形：光滑复曲面 $X$ 中两个除子 $D_1$、$D_2$ 的相交。它们对应
线丛 $L_i$ 与零点轨迹为 $D_i$ 的截面 $s_i\in H^0(L_i)$。$D_i$ 在 $H_2(X)$ 中有
同调类，可以求相交数；等价地，可以考察 $c_1(L_i)\in H^2(X)$。用层论的话说，
第一种做法是把结构层作张量积：若 $D_1$、$D_2$ 横截相交，则
$\OO_{D_1\cap D_2}=\OO_{D_1}\otimes\OO_{D_2}$。第二种做法是用消解
\eqref{eq:div} 表示 $\OO_{D_1}$，再把它与 $\OO_{D_2}$ 作张量积。下表把两种
做法并排放在一起：

\begin{table}[htbp]
\centering\small
\renewcommand{\arraystretch}{1.5}
\begin{tabular}{@{}p{2.4cm}p{4.6cm}p{6.8cm}@{}}
\toprule
数据 & 除子 $D_i$ & 截面 $s_i\in H^0(L_i)$，$s_i^{-1}(0)=D_i$\\
\midrule
代数几何 & $\OO_{D_i}$ & $\{L_i^{-1}\xrightarrow{s_i}\OO_X\}$\\
拓扑 & $[D_i]\in H_2(X)$ & $c_1(L_i)\in H^2(X)$\\
相交 & $[D_1]\cdot[D_2]$ &
$c_1(L_1)\cup c_1(L_2)=\langle c_1(L_1),D_2\rangle=c_1(L_1\res_{D_2})$\\
\midrule
横截情形 & $\OO_{D_1}\otimes\OO_{D_2}=\OO_{D_1\cap D_2}$ &
$c_1(L_1\res_{D_2})=(s_1\res_{D_2})^{-1}(0)$：把 $L_1$ 限制到 $D_2$，
取截面 $s_1\res_{D_2}$ 的零点，即取 $\{L_1^{-1}\res_{D_2}\xrightarrow{s_1\res_{D_2}}
\OO_{D_2}\}$ 的余核\\
\midrule
非横截情形 & $\OO_D\otimes\OO_D=\OO_D$ &
$\{L^{-1}\res_D\xrightarrow{s\res_D=0}\OO_D\}$：$D$ 上的线丛 $L$ 仍然在，
只是它的截面成了 $0$\\
（例如 $D_1=D_2=D$） & & \\
\bottomrule
\end{tabular}
\caption{除子、截面与相交数：代数几何与拓扑的对照（原表，无题注）}
\label{tab:intersection}
\end{table}

于是把 $L_1$ 限制到 $D_2$，相当于把 $\{L_1^{-1}\to\OO_X\}$ 与 $\OO_{D_2}$ 作
张量积；在横截情形，这个张量积的余核恰好是 $\OO_{D_1\cap D_2}$。因此，从右列
走到左列的办法是取余核；这对应一个事实：与 $\OO_{D_2}$ 作张量积是\emph{右正合}
的——把正合列 \eqref{eq:div} 与 $\OO_{D_2}$ 作张量积，得到的序列末三项仍然正合
\kanwuRef{2}。

这张表说明：把 $\OO_{D_1}$ 与 $\OO_{D_2}$ 作张量积，在 $D_1$、$D_2$ 横截时给出
正确答案，在非横截时则不然（$\OO_D$ 是 $D\cap D$ 的结构层，但不是 $D$ 与自身的
正确拓扑相交）。反过来，若把 $\{L_1^{-1}\to\OO_X\}$ 与 $\OO_{D_2}$ 作张量积，
则总能得到正确答案：取所得复形的余核，得到的仍是原来的
$\OO_{D_1}\otimes\OO_{D_2}$，但在非横截情形还多出了一些东西——线丛
$L_1\res_{D_2}$ 仍在，于是相交信息仍在（只需取它的一次陈类）；把可能在 $D_2$
的某些分支上恒为零的截面 $L_1^{-1}\res_{D_2}\xrightarrow{s_1\res_{D_2}}\OO_{D_2}$
换成横截的截面，就得到正确的相交数。更一般地，取 $L_1\res_{D_2}$ 在 $D_2$ 上的
除子。若能把 $D_1$ 移动成与 $D_2$ 横截，这给出同样的答案；若 $L_1\res_{D_2}$
没有截面，这一步做不到，但我们的方法仍然给出正确的相交积，而且是在概形论相交
$D_1\cap D_2$ \emph{内部}给出的。

教训是：
$$
\OO_D\otimes\OO_D\qquad\text{“本应是”}\qquad
\{L^{-1}\xrightarrow{s}\OO_X\}\otimes\OO_D
=\{L^{-1}\res_D\xrightarrow{0}\OO_D\}.
$$
后者的余核仍是 $\OO_D\otimes\OO_D=\OO_D$，但此时（因为相交不横截）还多出一个核，
即 $L^{-1}\res_D$，相交信息全在里面。这两个 $D$ 上的层的一次陈类之差，
恰好就是自相交数 $D\cdot D$。

所以说，正如 §\ref{sec:cx} 中单纯复形情形的 $H_*$ 的对偶“本应”是对链复形而非
$H_*$ 取 $\Hom(-,\Z)$，这里应当把 $\otimes\OO_D$ 作用到局部自由（因而性质更好）
的层组成的复形上，而不是作用到 $\OO_D$ 上。这是导出函子的原型，§\ref{sec:df}
会系统地处理。上面的核 $L^{-1}\res_D$ 就是 $\otimes\OO_D$ 的第一个导出函子
$\Tor_1(\OO_D,\OO_D)$\yizhuRef{9}。

于是我们看到一个例子，说明考察层的复形而不仅是单个（上同调）层是有好处的。
在作用了 $\otimes\OO_{D_2}$ 之后，我们得到的是一个真正的复形，上同调可能分布在
不止一个次数上，也就是说它不再是单个层的消解。但此时又不该马上取上同调，
因为例如当 $X$ 维数更高时，我们可能还要与别的圈相交，也就是把这个复形与某个
$\OO_{D_3}$ 作张量积。说到底，复形好，上同调坏。

既然已经为考察\emph{所有}复形（而不只是层的消解）找到了动机，下面就着手与它们
打交道，导出函子的事暂且搁下，留到 §\ref{sec:df}。把层换成以它为同调的复形
（即消解），又会碰到前面提过的问题：怎样函子化地挑一个消解？拓扑里同样有
对应的问题。

\begin{yizhu}[为什么核就是 $\Tor_1$]
两层 $A$、$B$ 的张量积的导出函子 $\Tor_i(A,B)$ 可用平坦（例如局部自由）消解
计算：取 $P\udot\to A$ 为平坦消解，则 $\Tor_i(A,B)=h_i(P\udot\otimes B)$。
在本文的例子里，$P\udot=\{L^{-1}\to\OO_X\}$、$B=\OO_D$，故
$\Tor_1(\OO_D,\OO_D)=h_1\{L^{-1}\res_D\xrightarrow{0}\OO_D\}=L^{-1}\res_D$，
$\Tor_0=\OO_D$。作者说的“第一个导出函子”，指的就是这个 $\Tor_1$。
[来源：原文此处；导出函子的定义见 §\ref{sec:df} 与 Gelfand--Manin 第 III 章；A 级]
\end{yizhu}

\begin{kanwu}[“末三个箭头”应为“末三项”]
原文说与 $\OO_{D_2}$ 作张量积后“the final three arrows are still exact”。
正合列 \eqref{eq:div} 只有三个箭头（四个对象）；张量积后，
$L^{-1}\res_{D_2}\to\OO_{D_2}\to\OO_{D_1\cap D_2}\to0$ 的“末三项”
（两条箭头）正合。原文“三个箭头”是“三项”的笔误。
译本正文按“末三项”处理，不改数学内容；详见附录 B 第 2 条。
\end{kanwu}

\begin{keytakeaways}
\item 同调群丢信息（Massey 积、同调相同而不同伦型），链复形才是“完整”的不变量；
  对偶要到复形层面做，才能把挠信息挪次数而不是丢掉。
\item 层论里的消解（除子正合列、Koszul 复形）把难缠的挠层换成局部自由层；
  代价是消解不典范。
\item 相交理论是导出的原型：$\OO_D\otimes\OO_D$ 只在横截时对，\\
  $\{L^{-1}\to\OO_X\}\otimes\OO_{D}$ 总对，多出的核就是 $\Tor_1$。
\item 一句话：复形好，（上）同调坏。
\end{keytakeaways}
